\thispagestyle{empty}
\pdfbookmark[0]{Edition and Sources}{translation-sources}
\begin{center}\textbf{Edition and Sources}\end{center}
\noindent This AI-translated edition presents the astronomical writings collected in Carl Friedrich Gauss’s \textit{Werke}, vol.~VI (Göttingen, 1874): papers, observations, excerpts from correspondence and reviews. The historical title pages and portrait accompany the text, together with Ernst Schering’s own editorial remarks. The translations follow the collected text, including its historical editorial selections and notes. References to pages of the original collected volume retain its 1874 pagination. Sources and translation parallels are identified in the accompanying notes. The title-page vignette is reproduced from the copy digitized by \href{https://archive.org/details/werkecarlf06gausrich/page/n8/mode/1up}{University of California Libraries}.
\par\medskip\noindent\textbf{Elliptic Elements of Pallas}
\par\smallskip
\noindent This AI translation covers all fifteen articles and the tables of Gauss’s \textit{Disquisitio de elementis ellipticis Palladis}. It follows the Latin text in \textit{Werke}, vol.~VI (Göttingen, 1874), including the paper’s historical title. The volume title pages and portrait also come from that edition. The paper was presented on 25 November 1810 and printed in volume~I of the \textit{Commentationes societatis regiae scientiarum Gottingensis recentiores} in 1811. \href{https://gdz.sub.uni-goettingen.de/id/PPN236007467}{Göttingen digitization}.
\par\smallskip\noindent Historical parallels include the German \href{https://www.digitale-sammlungen.de/de/view/bsb10538613?page=463}{editorial extract of November 1811}, the German version of parts of Articles~10--15 edited and translated by Anton Börsch and Paul Simon (Berlin, 1887, pp.~118--128), and Hale F. Trotter’s \href{https://www.york.ac.uk/depts/maths/histstat/gausspallas.htm}{English excerpt} from Joseph Bertrand’s French, in Princeton Technical Report No.~5 (1957). The Börsch and Simon collection has a foreword by Helmert. German orbital terminology was also compared with E. F. M. Jochim’s \textit{Satellitenbewegung}, vol.~III, version of 9 February 2022, pp.~60 and 63.
\par\smallskip\noindent Translation and typesetting produced by AI.
\clearpage
